\subsection{Best Practices}
\begin{tcolorbox}[colback=green!25]
    Compose a list of the best practices of energy usage and/or manufacturing activities in the plant by deleting inapplicable entries from the list below. 
\end{tcolorbox}
\noindent The plant operates its facility with several best energy efficiency practices: 
    \begin{itemize}\itemsep -2pt
        \item \emph{LED lighting system} has been installed in part of the facility.
        \item \emph{Proper rate schedules} for utilities.
        \item \emph{Roof renovations} was recently performed.
        \item \emph{Refrigerated dryers} are installed after each compressor.
        \item \emph{Noise level awareness} for loud areas and signs posted requiring hearing protection.
        \item \emph{Statistical quality control} management system overview of each machine.
        \item \emph{BMS system} managing all major equipment in the facility.
        \item \emph{Backup generators} for safety/emergency.
        \item \emph{VFDs} installed for many motor-driven equipment (pumps, blower fans, etc.).
        \item \emph{Notched motor belts} are used for belt-driven equipment (e.g., ventilation fans).
        \item \emph{Continued commitment to energy efficiency} for retrofitting and seeking efficiency measures to reduce energy expenditures.
        \item \emph{Occupancy sensors} installed in low-traffic areas to reduce unnecessary lighting.
        \item \emph{Programmable thermostats} to optimize HVAC scheduling during occupied and unoccupied periods.
        \item \emph{High-efficiency motors} (NEMA Premium or IE3/IE4) used for critical equipment.
        \item \emph{Power factor correction} capacitors installed to minimize reactive power charges.
        \item \emph{Compressed air leak detection program} with regular audits to identify and repair leaks.
        \item \emph{Receiver tanks} properly sized to reduce compressor cycling and short-loading.
        \item \emph{Insulated steam and hot water pipes} to minimize thermal losses in distribution systems.
        \item \emph{Steam trap maintenance program} with regular inspections and replacements.
        \item \emph{Economizer controls} on boilers to preheat feedwater using flue gas heat recovery.
        \item \emph{Destratification fans} installed in high-bay areas to recirculate warm air during heating season.
        \item \emph{High-speed doors} or strip curtains at loading docks to reduce infiltration losses.
        \item \emph{Demand-controlled ventilation} adjusting outside air intake based on occupancy or CO2 levels.
        \item \emph{Premium air filters} with appropriate MERV ratings balanced for air quality and pressure drop.
        \item \emph{Reflective or cool roofing materials} to reduce solar heat gain during cooling season.
        \item \emph{Preventive maintenance program} with scheduled equipment inspections and servicing.
        \item \emph{Energy submetering} installed on major systems to track consumption and identify anomalies.
        \item \emph{Waste heat recovery} from compressors, ovens, or process equipment for space or water heating.
        \item \emph{Natural daylighting} through skylights or clerestory windows in production areas.
        \item \emph{Automatic shutoff controls} for equipment during breaks, shift changes, or weekends.
        \item \emph{Condensate return system} to recycle hot condensate back to the boiler feedwater tank.
        \item \emph{Energy awareness training} provided to employees on conservation practices.
    \end{itemize}